\chapter[CƠ SỞ LÝ THUYẾT]{Cơ sở lý thuyết}

Chương này trình bày các khái niệm và phương pháp nền tảng phục vụ xây dựng mô hình ở chương sau, theo trình tự: từ khái quát học máy, xử lý ngôn ngữ tự nhiên, các kỹ thuật trích xuất đặc trưng văn bản, các mô hình học máy truyền thống, thuật toán láng giềng gần nhất (KNN) --- thuật toán chính của đề tài --- và các độ đo đánh giá mô hình.

\section{Học máy}

\subsection{Khái niệm học máy}

\textbf{Học máy} (Machine Learning --- ML) là nhánh của trí tuệ nhân tạo nghiên cứu các thuật toán cho phép máy tính \textit{học} từ dữ liệu thay vì được lập trình trực tiếp bằng quy tắc cứng. Như Arthur Samuel mô tả, học máy là lĩnh vực cho phép máy tính có khả năng học mà không cần được lập trình rõ ràng. Theo tiếp cận học máy của Tom Mitchell, một chương trình học từ \textbf{tập kinh nghiệm} $E$ đối với \textbf{tác vụ} $T$ và \textbf{độ đo} $P$; nếu hiệu năng thực hiện $T$ đo bằng $P$ được cải thiện theo $E$ thì chương trình đó được coi là ``học'' \cite{tiep}.

Các bài toán học máy được chia thành ba nhóm chính:

\begin{itemize}
    \item \textbf{Học có giám sát} (supervised learning): học từ dữ liệu có gán nhãn, ví dụ phân loại (classification) và hồi quy (regression). Bài toán phát hiện tin giả là bài toán học có giám sát dạng phân loại nhị phân.
    \item \textbf{Học không giám sát} (unsupervised learning): học từ dữ liệu không có nhãn, ví dụ phân cụm (clustering), giảm chiều (dimensionality reduction).
    \item \textbf{Học tăng cường} (reinforcement learning): tác nhân học hành vi qua tương tác với môi trường bằng phần thưởng/trị phạt.
\end{itemize}

\subsection{Học có giám sát và bài toán phân loại}

Trong học có giám sát, thuật toán nhận vào một tập huấn luyện gồm các cặp (đặc trưng, nhãn) $(x_i, y_i)$ và xây dựng một hàm dự đoán $f$ sao cho $f(x) \approx y$ trên dữ liệu mới. Khi nhãn là biến rời rạc (ví dụ ``tin thật''/``tin giả''), bài toán gọi là \textbf{phân loại}; khi nhãn là biến liên tục, bài toán gọi là \textbf{hồi quy}.

Các thuật toán phân loại có thể chia theo cách ``lưu trữ tri thức'' thành hai nhóm đáng chú ý:

\begin{itemize}
    \item \textbf{Học nhiệt tình} (eager learning): trong giai đoạn huấn luyện, thuật toán tổng kết toàn bộ dữ liệu thành một mô hình tham số compact --- ví dụ bộ trọng số của mô hình tuyến tính --- rồi loại bỏ dữ liệu gốc; khi dự đoán chỉ cần mô hình này.
    \item \textbf{Học dựa trên mẫu} (instance-based learning): thuật toán không xây dựng mô hình tổng quát mà \textit{lưu giữ nguyên} các mẫu huấn luyện; khi cần dự đoán một mẫu mới, thuật toán so sánh trực tiếp mẫu đó với các mẫu đã lưu. Thuật toán láng giềng gần nhất (KNN) là đại diện tiêu biểu của nhóm này, và được gọi là thuật toán \textbf{lười} (lazy learning) vì công sức tính toán bị dồn sang giai đoạn dự đoán thay vì giai đoạn huấn luyện.
\end{itemize}

\subsection{Quy trình xây dựng một bài toán học máy}

Quy trình chuẩn cho một bài toán học máy có giám sát được thực hiện trong đề tài gồm các bước:

\begin{enumerate}
    \item \textbf{Xác định bài toán}: xác định rõ đầu vào (văn bản tin tức), đầu ra (nhãn thật/giả) và tiêu chí đánh giá thành công (Accuracy, F1-score\ldots).
    \item \textbf{Thu thập dữ liệu}: tìm bộ dữ liệu gán nhãn chất lượng, đủ lớn và có nguồn gốc rõ ràng (bộ ISOT \cite{dataset}).
    \item \textbf{Khảo sát dữ liệu (EDA)}: thống kê số lượng mẫu, phân bố nhãn, độ dài văn bản, kiểm tra giá trị thiếu, trùng lặp.
    \item \textbf{Tiền xử lý}: làm sạch văn bản, chuẩn hóa, loại bỏ nhiễu.
    \item \textbf{Trích xuất đặc trưng / vector hóa}: chuyển văn bản thành dạng số mà thuật toán học được (TF-IDF).
    \item \textbf{Chia tập dữ liệu}: chia ngẫu nhiên thành tập huấn luyện (train) để học tham số và tập kiểm tra (test) để đánh giá khách quan. Trong đề tài dùng tỷ lệ 75\%/25\%. Với các bài toán lớn, thường chia thêm tập phát triển (validation) để tinh chỉnh siêu tham số hoặc dùng \textbf{kiểm chứng chéo k-lớp} (k-fold cross-validation): chia dữ liệu thành $k$ phần, luân phiên lấy mỗi phần làm tập kiểm tra và lấy trung bình kết quả.
    \item \textbf{Huấn luyện}: chạy thuật toán tối ưu để tìm tham số tốt nhất trên tập huấn luyện. Với KNN, bước này chỉ là ghi nhớ tập dữ liệu (chi phí gần như bằng 0), mọi tính toán dồn sang bước dự đoán.
    \item \textbf{Đánh giá và phân tích lỗi}: đo các độ đo đánh giá trên tập kiểm tra chưa từng thấy, phân tích các mẫu bị phân loại sai.
    \item \textbf{Triển khai}: đóng gói mô hình vào ứng dụng thực tế.
\end{enumerate}

\section{Xử lý ngôn ngữ tự nhiên}

\textbf{Xử lý ngôn ngữ tự nhiên} (Natural Language Processing --- NLP) là lĩnh vực giao thoa giữa ngôn ngữ học, trí tuệ nhân tạo và khoa học máy tính, nghiên cứu cách máy tính có thể hiểu, phân tích và tạo ra ngôn ngữ tự nhiên của con người. Phân loại văn bản (text classification) là một trong những bài toán cơ bản nhất của NLP, bao gồm phân loại cảm xúc, lọc thư rác, phân loại chủ đề tin tức và cả phát hiện tin giả.

Khác với dữ liệu bảng truyền thống, văn bản có đặc điểm: độ dài thay đổi, thứ tự từ mang ý nghĩa, từ vựng khổng lồ, và cùng một nội dung có thể diễn đạt vô số cách. Do đó trước khi đưa vào mô hình, văn bản cần được chuẩn hóa qua các bước xử lý cơ bản:

\begin{itemize}
    \item \textbf{Chuẩn hóa chữ}: chuyển toàn bộ văn bản về chữ thường (lowercase) để coi ``News'' và ``news'' là cùng một từ; loại bỏ dấu câu, ký tự đặc biệt, đường dẫn, số và thẻ HTML --- các thành phần ít mang thông tin phân loại.
    \item \textbf{Tách từ} (tokenization): phân tách văn bản thành các đơn vị nhỏ hơn (từ hoặc ký tự), là đơn vị cơ bản để xây dựng từ điển và vector hóa.
    \item \textbf{Loại bỏ từ dừng} (stopwords removal): các từ xuất hiện rất thường xuyên nhưng ít mang thông tin phân loại như ``a'', ``the'', ``is'', ``and''\ldots được loại bỏ theo bộ từ dừng chuẩn của từng ngôn ngữ (NLTK).
    \item \textbf{Rút gọn từ} (stemming/lemmatization): đưa các biến thể của từ về dạng gốc. \textit{Stemming} dùng thuật toán cắt đuôi nhanh (Porter Stemmer) có thể cho kết quả không phải từ đúng chính tả (``studies'' $\rightarrow$ ``studi''); \textit{lemmatization} tra cứu từ điển cho ra từ nguyên chuẩn (``studies'' $\rightarrow$ ``study'') nhưng chậm hơn. Với bài toán phân loại ở quy mô lớn, stemming thường đủ hiệu quả.
\end{itemize}

\textbf{Mô hình n-gram} là khái niệm quan trọng liên quan: một n-gram là dãy $n$ từ liên tiếp. Unigram là từ đơn, bigram là cặp từ, trigram là cụm 3 từ. Việc tính đặc trưng theo n-gram cho phép giữ lại một phần ngữ nghĩa bị mất bởi mô hình túi từ, ví dụ bigram ``not good'' mang ý nghĩa trái ngược với hai unigram ``not'' và ``good'' riêng lẻ.

\section{Trích xuất đặc trưng văn bản}

\subsection{Mô hình túi từ}

\textbf{Mô hình túi từ} (Bag of Words) biểu diễn văn bản bằng vector tần suất xuất hiện của các từ trong từ điển, \textit{bỏ qua hoàn toàn thứ tự} của chúng. Ví dụ với từ điển 4 từ \{``tin'', ``giả'', ``lan'', ``truyền''\}, câu ``tin giả lan truyền'' được biểu diễn [1, 1, 1, 1]. Mô hình đơn giản, dễ triển khai nhưng: từ điển khổng lồ gây vector thưa (sparse) kích thước lớn; mất hoàn toàn thông tin thứ tự; các từ hiếm gây nhiễu.

\subsection{TF-IDF}
\label{sec:tfidf}

\textbf{TF-IDF} (Term Frequency -- Inverse Document Frequency) là kỹ thuật đánh trọng số của từ trong văn bản, được tính theo công thức:

\begin{equation}
    \mathrm{TFIDF}(t, d) = \mathrm{TF}(t, d) \times \mathrm{IDF}(t)
    \label{eq:tfidf}
\end{equation}

với $\mathrm{TF}(t, d)$ là tần suất của từ $t$ trong văn bản $d$ (có thể chuẩn hóa theo độ dài văn bản) và

\begin{equation}
    \mathrm{IDF}(t) = \log \frac{N}{1 + n_t}
    \label{eq:idf}
\end{equation}

trong đó $N$ là tổng số văn bản trong kho dữ liệu và $n_t$ là số văn bản chứa từ $t$. Ý nghĩa của công thức: từ xuất hiện ở nhiều văn bản (ít mang tính phân biệt, ví dụ từ dừng) có IDF nhỏ; từ xuất hiện tập trung ở một số văn bản (đặc trưng cho chủ đề) có IDF lớn. Nhân với TF, những từ \textit{vừa đặc trưng cho văn bản vừa hiếm trong kho} nhận trọng số cao nhất.

\textbf{Ví dụ minh họa.} Giả sử kho dữ liệu có $N = 1.000$ bài tin tức; từ ``bầu'' xuất hiện 900 bài (IDF $\approx \log(1000/901) \approx 0{,}05$ --- gần như vô nghĩa) trong khi từ ``đàn áp'' chỉ xuất hiện 50 bài (IDF $\approx \log(1000/51) \approx 1{,}3$ --- mang tính phân biệt cao). Với cùng tần suất xuất hiện trong một bài, trọng số TF-IDF của ``đàn áp'' sẽ gấp khoảng 25 lần ``bầu''.

\textbf{Chuẩn hóa L2 của vector TF-IDF.} Sau khi tính trọng số, mỗi vector văn bản thường được \textit{chuẩn hóa L2}: chia toàn bộ thành phần cho độ dài Euclidean của vector, khi đó $\lVert x \rVert = 1$. Chuẩn hóa này đưa mọi bài viết về cùng một ``thang đo'' bất kể bài dài hay ngắn, tránh việc bài dài có tổng trọng số lớn hơn một cách ``không công bằng''. Đối với thuật toán KNN, chuẩn hóa L2 còn có ý nghĩa đặc biệt: với hai vector chuẩn hóa $x, y$, khoảng cách Euclidean và khoảng cách cosine liên hệ với nhau qua đẳng thức
\begin{equation}
    \lVert x - y \rVert^2 = \lVert x \rVert^2 + \lVert y \rVert^2 - 2\,x^{\top}y = 2\left(1 - \cos(x, y)\right),
    \label{eq:l2-cos}
\end{equation}
tức hai phép đo chỉ khác nhau một biến đổi đơn điệu trên cùng dữ liệu --- đây là lý do kinh điển khiến cosine là phép đo được ưa dùng cho văn bản (phân tích chi tiết ở mục \ref{sec:knn-khoang-cach}).

TF-IDF kết hợp với phân tích n-gram và mô hình học máy truyền thống (Naive Bayes, Logistic Regression, Passive Aggressive, KNN) là phương pháp kinh điển cho phân loại văn bản và vẫn là baseline mạnh tới ngày nay.

\section{Một số mô hình học máy truyền thống}

Các mô hình học máy truyền thống kết hợp đặc trưng TF-IDF là cơ sở so sánh quan trọng trong đề tài. Dưới đây là nguyên lý hoạt động của ba mô hình tiêu biểu được dùng làm tham chiếu.

\subsection{Naive Bayes}

\textbf{Naive Bayes} là mô hình phân loại dựa trên định lý Bayes với giả định ``ngây thơ'' (naive): các đặc trưng độc lập với nhau khi biết lớp của mẫu. Với bài toán phân loại tin thật/giả, mẫu $x$ được gán vào lớp có xác suất sau lớn nhất:

\begin{equation}
    \hat{y} = \arg\max_{c \in \{0, 1\}} P(c) \prod_{i=1}^{n} P(x_i \mid c)
    \label{eq:nb}
\end{equation}

trong đó $P(c)$ là xác suất tiền nghiệm của lớp $c$ (ước lượng từ tỷ lệ các bài thật/giả trong tập huấn luyện) và $P(x_i \mid c)$ là xác suất từ $x_i$ xuất hiện trong lớp $c$. Với đặc trưng TF-IDF là giá trị liên tục, thường dùng biến thể \textbf{Multinomial Naive Bayes} coi mỗi giá trị đặc trưng như số lần xuất hiện của từ. Ưu điểm: huấn luyện cực nhanh, ít tham số, hiệu quả với dữ liệu thưa; nhược điểm: giả định độc lập giữa các từ rõ ràng sai (ngôn ngữ có cấu trúc), nên chất lượng bị giới hạn.

\subsection{Logistic Regression}

\textbf{Hồi quy Logistic} mô hình hóa xác suất mẫu thuộc lớp dương qua hàm sigmoid áp lên một tổ hợp tuyến tính của đặc trưng:

\begin{equation}
    P(y = 1 \mid x) = \sigma(w^{\top} x + b) = \frac{1}{1 + e^{-(w^{\top}x + b)}}
    \label{eq:lr}
\end{equation}

với $w$ là vector trọng số của các đặc trưng TF-IDF và $b$ là thiên vị. Trọng số được học bằng tối ưu hóa \textit{maximum likelihood} thường kết hợp chuẩn hóa $L_2$ để chống quá khớp. Ý nghĩa của trọng số $w_i$ trực quan: từ có trọng số dương lớn (``shocking'', ``miracle'') đẩy dự đoán về phía tin giả, từ có trọng số âm lớn (``Reuters'' --- dấu hiệu nguồn chính thống) đẩy về phía tin thật. Mô hình đơn giản, dễ diễn giải, và trên dữ liệu văn bản với đặc trưng TF-IDF thường đạt chất lượng rất cao --- tiêu biểu cho nhóm phương pháp tuyến tính.

\subsection{Passive Aggressive}

\textbf{Passive Aggressive} là mô hình học trực tuyến (online learning) cho phân loại: với mỗi mẫu mới, mô hình chỉ ``chủ động'' (aggressive) cập nhật trọng số khi phân loại sai hoặc dự đoán với biên độ không đủ lớn, ngược lại giữ nguyên (passive). Cập nhật có dạng:

\begin{equation}
    w_{t+1} = w_t + \lambda_t \, y_t \, x_t
    \label{eq:pa}
\end{equation}

trong đó $\lambda_t$ được chọn là mức cập nhật tối thiểu đủ để sửa mẫu hiện tại. Thuật toán phù hợp với luồng tin tức liên tục thay đổi, và là mô hình tham chiếu được sử dụng trong nghiên cứu gốc của bộ dữ liệu ISOT \cite{ahmed2017}.

\section{Thuật toán láng giềng gần nhất (KNN)}

\subsection{Nguyên lý hoạt động}

\textbf{Thuật toán láng giềng gần nhất} (K-Nearest Neighbors --- KNN) do Cover và Hart đề xuất \cite{coverhart} là một trong những thuật toán phân loại đơn giản và trực quan nhất của học máy, dựa trên nguyên lý ``\textit{hãy cho tôi biết bạn sống cùng ai, tôi sẽ nói bạn là ai}'': một mẫu mới được gán nhãn theo \textbf{số đông} trong $k$ mẫu đã biết \textit{gần nhất} (tương đồng nhất) với nó trong không gian đặc trưng.

Khác với các mô hình ở mục trước, KNN là thuật toán \textbf{lười} (lazy learning): giai đoạn ``huấn luyện'' chỉ gồm việc lưu trữ toàn bộ tập huấn luyện; toàn bộ công sức tính toán dồn vào lúc dự đoán. Quy trình dự đoán một bài tin tức mới $x$ gồm các bước:

\begin{enumerate}
    \item Tính khoảng cách từ $x$ tới \textbf{tất cả} các mẫu trong tập huấn luyện.
    \item Chọn ra $k$ mẫu có khoảng cách nhỏ nhất --- gọi là $k$ \textbf{láng giềng gần nhất}, ký hiệu tập $N_k(x)$.
    \item Gán nhãn theo \textbf{bỏ phiếu đa số} (hoặc bỏ phiếu có trọng số, xem mục \ref{sec:knn-trong-so}):
    \begin{equation}
        \hat{y} = \arg\max_{c \in \{0, 1\}} \sum_{i \in N_k(x)} w_i \cdot \mathbb{1}\left[y_i = c\right],
        \label{eq:knn-vote}
    \end{equation}
    trong đó $\mathbb{1}[\cdot]$ là hàm chỉ báo (bằng 1 khi điều kiện đúng) và $w_i$ là trọng số của láng giềng $i$.
\end{enumerate}

Sơ đồ \ref{sodo:knn-nguyen-ly} minh họa nguyên lý trên một ví dụ phẳng hai chiều: bài viết cần dự đoán (dấu ``?'') nằm gần 4 bài tin giả và 1 bài tin thật, nên được dự đoán là tin giả.

\begin{sodo}[H]
    \centering
    \begin{tikzpicture}[>=stealth, font=\small]
        % Cac diem tin that (xanh duong)
        \foreach \x/\y in {1/4.2, 2/3.4, 1.6/5.2, 3/5, 2.4/4.6, 0.8/5.6, 1/6.2, 2.2/6} {
            \draw[fill=blue!35] (\x,\y) circle (0.15);
        }
        % Cac diem tin gia (do)
        \foreach \x/\y in {6.6/2.6, 6.4/2, 6.8/3.6, 6.2/5, 5/5.4} {
            \draw[fill=red!35] (\x,\y) circle (0.15);
        }
        % 5 lang gieng gan nhat cua diem can du doan
        \coordinate (n1) at (4.4,3);
        \coordinate (n2) at (5,2.2);
        \coordinate (n3) at (5.4,3.8);
        \coordinate (n4) at (5.8,4.6);
        \coordinate (n5) at (3.4,4);
        \foreach \p in {n1,n2,n3,n4} {\draw[fill=red!35] (\p) circle (0.15);}
        \draw[fill=blue!35] (n5) circle (0.15);
        % Diem can du doan
        \node[draw, thick, fill=yellow!60, circle, minimum size=0.55cm] (q) at (4.5,3.3) {\textbf{?}};
        % Vong tron khoang cach k lang gieng
        \draw[dashed, thick] (q) circle (1.9);
        % Mui ten tu diem can du doan toi 5 lang gieng
        \foreach \p in {n1,n2,n3,n4,n5} {\draw[->, thick] (q) -- (\p);}
        % Chu thich
        \draw[fill=blue!35] (-1.4,0.3) circle (0.13);
        \node[anchor=west, font=\footnotesize] at (-1.2,0.3) {Tin thật (đã gán nhãn)};
        \draw[fill=red!35] (1.6,0.3) circle (0.13);
        \node[anchor=west, font=\footnotesize] at (1.8,0.3) {Tin giả (đã gán nhãn)};
        \node[draw, thick, fill=yellow!60, circle, minimum size=0.42cm, inner sep=0pt] at (4.3,0.3) {\footnotesize\textbf{?}};
        \node[anchor=west, font=\footnotesize] at (4.6,0.3) {Bài viết cần dự đoán};
    \end{tikzpicture}
    \SodoCaption{Nguyên lý KNN: dựa trên 5 láng giềng gần nhất (4 tin giả, 1 tin thật) bài viết được dự đoán là tin giả}{Nhóm tự xây dựng theo \cite{coverhart}}
    \label{sodo:knn-nguyen-ly}
\end{sodo}

\subsection{Các phép đo khoảng cách}
\label{sec:knn-khoang-cach}

Kết quả của KNN phụ thuộc trực tiếp vào cách đo \textbf{độ tương đồng} giữa hai mẫu. Ba phép đo khoảng cách phổ biến nhất:

\begin{itemize}
    \item \textbf{Khoảng cách Euclidean} --- độ dài đoạn thẳng nối hai điểm trong không gian đặc trưng:
    \begin{equation}
        d(x, y) = \sqrt{\sum_{j=1}^{n} \left(x_j - y_j\right)^2}
        \label{eq:euclid}
    \end{equation}
    \item \textbf{Khoảng cách Manhattan} --- tổng độ chênh lệch theo từng thành phần (đi theo ``ô phố vuông''):
    \begin{equation}
        d(x, y) = \sum_{j=1}^{n} \left| x_j - y_j \right|
        \label{eq:manhattan}
    \end{equation}
    \item \textbf{Khoảng cách cosine} --- đo độ lệch \textit{góc} giữa hai vector, bất kể độ dài vector:
    \begin{equation}
        d(x, y) = 1 - \cos(x, y) = 1 - \frac{x^{\top} y}{\lVert x \rVert \, \lVert y \rVert}
        \label{eq:cos}
    \end{equation}
\end{itemize}

Đối với dữ liệu văn bản, \textbf{cosine} là lựa chọn được ưa dùng: hai bài viết cùng chủ đề nhưng một bài dài một bài ngắn có thể có vector TF-IDF ``chỉ về cùng một hướng'' --- cosine coi chúng tương đồng cao, trong khi Euclidean bị chi phối bởi độ lớn các thành phần. Như đã chỉ ra ở đẳng thức \eqref{eq:l2-cos}, khi các vector đã được chuẩn hóa L2, khoảng cách Euclidean và khoảng cách cosine tạo ra \textit{cùng một thứ tự} láng giềng trên cùng dữ liệu; khác biệt then chốt nằm ở cách thư viện cài đặt phép tính trên ma trận thưa siêu chiều, sẽ được phân tích kỹ trong thí nghiệm E4 ở chương 3.

\subsection{Lựa chọn số láng giềng k}
\label{sec:knn-k}

Tham số quan trọng nhất của KNN là số láng giềng $k$:

\begin{itemize}
    \item \textbf{$k$ quá nhỏ} (ví dụ $k = 1$): mô hình chỉ dựa vào một láng giềng duy nhất, rất nhạy với nhiễu --- một bài viết dị biệt (ví dụ tin thật viết theo văn phong giật gân) có thể ``kéo'' cả dự đoán về phía sai. Ranh giới quyết định trở nên răng cưa, \textit{phương sai} cao.
    \item \textbf{$k$ quá lớn}: dự đoán bị ``làm mượt'' bởi quá nhiều láng giềng, kể cả những láng giềng đã khá xa và ít liên quan, khiến ranh giới quá thô, \textit{độ lệch} cao.
    \item Với bài toán \textbf{phân loại nhị phân}, $k$ nên chọn \textbf{số lẻ} để tránh hiện tượng hòa phiếu khi dùng bỏ phiếu đa số.
\end{itemize}

Trong thực hành, $k$ được xác định bằng kiểm chứng chéo trên tập phát triển; trong khuôn khổ đề tài, $k$ được khảo sát qua thí nghiệm E3 (chương 3) với $k \in \{1, 3, 5, 7, 9, 11\}$ trên cùng cách chia dữ liệu, và cấu hình chuẩn chọn $k = 5$ --- giá trị cân bằng phổ biến trong tài liệu \cite{sklearn}.

\subsection{KNN có trọng số}
\label{sec:knn-trong-so}

Bỏ phiếu đa số thuần túy (\textit{uniform}) coi mọi láng giềng ngang hàng, dù láng giềng đó ở rất gần hay chỉ ``vừa đủ'' nằm trong top $k$. Phiên bản \textbf{KNN có trọng số theo khoảng cách} (\textit{distance}) khắc phục điều này bằng cách gán cho láng giềng thứ $i$ trọng số tỉ lệ nghịch với khoảng cách:

\begin{equation}
    w_i = \frac{1}{d(x, x_i)}
    \label{eq:distance-weight}
\end{equation}

Nhờ đó, một láng giềng ở rất gần (văn bản gần như trùng lặp về vector) có thể ``áp đảo'' nhiều láng giềng ở xa cộng lại. Trọng số theo khoảng cách còn cho ra \textbf{xác suất dự đoán mềm}: tỷ lệ trọng số của từng lớp trong $k$ láng giềng được dùng làm xác suất, cho phép vẽ đường cong ROC và thay đổi ngưỡng quyết định (mục độ đo ở cuối chương) thay vì chỉ nhận nhãn cứng.

\subsection{Ưu nhược điểm và độ phức tạp}

\textbf{Ưu điểm của KNN:}
\begin{itemize}
    \item \textbf{Đơn giản, dễ cài đặt, dễ diễn giải}: kết quả dự đoán có thể ``giải thích'' ngay bằng danh sách $k$ bài viết gần nhất và nhãn của chúng --- đặc điểm quý giá cho bài toán phát hiện tin giả, nơi con người cần kiểm tra lại quyết định của máy.
    \item \textbf{Không cần giai đoạn huấn luyện}: thêm dữ liệu mới (bài tin mới được xác thực) chỉ là thêm một mẫu vào kho lưu trữ, không phải huấn luyện lại mô hình.
    \item \textbf{Ranh giới quyết định phi tuyến tự nhiên}: KNN không giả định bất kỳ dạng hàm nào (tuyến tính, đa thức\ldots), phù hợp dữ liệu có cấu trúc phức tạp.
\end{itemize}

\textbf{Nhược điểm của KNN:}
\begin{itemize}
    \item \textbf{Chi phí dự đoán lớn}: để dự đoán một mẫu, phải tính khoảng cách tới toàn bộ $n$ mẫu huấn luyện, độ phức tạp $O(n \cdot d)$ với $d$ là số chiều đặc trưng. Với $n = 29.328$ bài huấn luyện và vector 50.000 chiều, dự đoán 9.777 bài kiểm tra mất hàng phút --- chậm hơn hàng nghìn lần các mô hình tuyến tính (đo cụ thể ở chương 3). Đây là điểm khác biệt căn bản với học nhiệt tình: mô hình tuyến tính trả ``tiền'' ở lúc huấn luyện, KNN trả ``tiền'' ở mỗi lần dự đoán.
    \item \textbf{Tốn bộ nhớ}: phải giữ nguyên toàn bộ tập huấn luyện.
    \item \textbf{Nhạy cảm với đặc trưng không liên quan và thước đo khoảng cách}: trên không gian chiều rất cao (lời nguyền chiều, curse of dimensionality), các khoảng cách có xu hướng đồng nhất và ý nghĩa ``gần nhau'' bị suy giảm; việc chọn đặc trưng và phép đo khoảng cách quyết định chất lượng mô hình.
\end{itemize}

\section{Các độ đo đánh giá mô hình}

Với bài toán phân loại nhị phân, chất lượng mô hình được đánh giá trên \textbf{ma trận nhầm lẫn} (confusion matrix) gồm 4 đại lượng:

\begin{itemize}
    \item \textbf{TP} (True Positive): tin giả được dự đoán đúng là giả.
    \item \textbf{TN} (True Negative): tin thật được dự đoán đúng là thật.
    \item \textbf{FP} (False Positive): tin thật bị dự đoán nhầm là giả.
    \item \textbf{FN} (False Negative): tin giả bị dự đoán nhầm là thật.
\end{itemize}

Từ đó tính được các độ đo:

\begin{itemize}
    \item \textbf{Độ chính xác} (Accuracy) --- tỷ lệ dự đoán đúng trên tổng số mẫu:
    \begin{equation}
        \mathrm{Accuracy} = \frac{TP + TN}{TP + TN + FP + FN}
    \end{equation}
    \item \textbf{Độ chuẩn xác} (Precision) --- trong các bài bị gắn cờ giả, bao nhiêu bài thực sự là giả:
    \begin{equation}
        \mathrm{Precision} = \frac{TP}{TP + FP}
    \end{equation}
    \item \textbf{Độ bao phủ} (Recall) --- trong các bài giả thực tế, bao nhiêu bài được bắt được:
    \begin{equation}
        \mathrm{Recall} = \frac{TP}{TP + FN}
    \end{equation}
    \item \textbf{Điểm F1} (F1-score) --- trung bình hài hòa của Precision và Recall, cân bằng hai chỉ tiêu trên:
    \begin{equation}
        F1 = 2 \cdot \frac{\mathrm{Precision} \times \mathrm{Recall}}{\mathrm{Precision} + \mathrm{Recall}}
    \end{equation}
\end{itemize}

\textbf{Lưu ý quan trọng khi chọn độ đo:} Accuracy có thể gây hiểu nhầm khi dữ liệu lệch lớp. Ví dụ, nếu 95\% bài viết là tin thật, một mô hình ``lười'' dự đoán mọi bài là thật vẫn đạt Accuracy 95\% nhưng Recall trên tin giả bằng 0 --- vô dụng trong thực tế. Trong phát hiện tin giả, \textbf{Recall} (bắt được càng nhiều tin giả càng tốt) và \textbf{Precision} (ít ``báng bổ'' các tin thật hợp lệ) đặc biệt quan trọng vì chi phí của hai loại sai khác nhau: bỏ sót tin giả cho phép tin độc hại lan truyền, còn gắn cờ nhầm tin thật làm giảm uy tín hệ thống.

Một công cụ phân tích khác là \textbf{đường cong ROC} (Receiver Operating Characteristic) vẽ tỷ lệ TP theo tỷ lệ FP khi thay đổi ngưỡng quyết định, và \textbf{diện tích dưới đường cong AUC} (Area Under Curve): AUC = 1 là mô hình hoàn hảo, AUC = 0{,}5 tương đương dự đoán ngẫu nhiên. AUC đo chất lượng phân tách của mô hình \textit{không phụ thuộc ngưỡng}.

Trong bộ dữ liệu ISOT, hai lớp gần cân bằng (54,2\%/45,8\%), nên Accuracy là chỉ tiêu so sánh chính trong báo cáo này, song song với F1-score và ma trận nhầm lẫn để phân tích chi tiết từng loại sai.
